\section{Demand (Pierre-Loup)}\label{sec:demand}

Demand allocation is the selection environment: it sets how fast a price advantage becomes
market share, and so whether a price war pays. The selection intensity $\chi$ may be the
most important parameter of the project. With strong selection a small price cut moves
customers fast, price wars pay and the unit-cost rule gains; with weak selection shares
barely move and firms with higher profit rates win by accumulating faster. Selection is weak
in the data (Dosi et al.\ 2015; Foster, Haltiwanger \& Syverson 2008), so results must be
reported at weak selection too.

The course's K+S formula uses a negative competitiveness; divided by a negative mean, the
update would send customers to the dearer firm. Version 1 uses the positive
competitiveness $E = 1/p$. Industry demand is insensitive to the industry price, so price
wars only move customers between firms.

\paragraph{Starting rule.}
\[
E_{i,t} = \frac{1}{p_{i,t}}, \qquad \bar E_t = \sum_j f_{j,t-1}\, E_{j,t}, \qquad
f_{i,t} = f_{i,t-1}\left(1 + \chi\, \frac{E_{i,t} - \bar E_t}{\bar E_t}\right), \quad \chi = 1.
\]
Shares cut below zero are set to zero and the others renormalised. A firm sells
\[
x_{i,t} = \min\{f_{i,t} D_t,\ \text{output it offers at } p_{i,t}\},
\]
where the output
offered is the capacity of the plants whose operating cost the price covers, up to
$u_{\max}\qbar_i$. The overflow goes to the cheapest firms with spare capacity (tied firms
share it in proportion to spare capacity); each firm's unmet demand and the demand nobody
can serve are recorded. Utilisation is $u_{i,t} = x_{i,t}/\qbar_i$. At $\chi = 1$, a small
firm 2\,\% cheaper than average grows its share by about 2\,\% a period.

The alternative \code{DEMAND\_FORM = "logit"} sets $f_{i,t} \propto f_{i,t-1}\, p_{i,t}^{-\eta}$,
which never produces negative shares; with $\eta = 1$ it responds like $\chi = 1$ to small
gaps.

\paragraph{Tests} (\code{tests/test\_demand.py}): shares non-negative and summing to one;
equal prices leave shares unchanged; a firm 2\,\% cheaper gains share at the rate $\chi$
predicts; overflow goes to others and no unit is sold twice; ties are shared pro rata;
swapping the type labels changes nothing.

\subsection*{Code}
\eqfile{competitiveness.py}
\eqfile{market_shares.py}
\eqfile{offered_capacity.py}
\marketfile{goods_market.py}
\schedulerpart{demand}{\code{model/scheduler.py}: step 2}
